% ---
% title: Rethinking the Design Space of RL for Diffusion Models — On the Importance of Likelihood Estimation Beyond Loss Design
% date: 2026-03-28
% tag: Diffusion
% excerpt: 本文系统分析 RL fine-tuning diffusion model 的三大设计维度，发现 ELBO-based likelihood estimation 是主导效率与性能的核心因素，而非 policy-gradient loss 的具体形式；最终方法在 90 GPU hours 内将 SD3.5-M 的 GenEval 从 0.24 提升至 0.95。
% ---
\section{背景与动机}

Diffusion/Flow model（DDPM、Stable Diffusion、SD3.5 等）已成为 text-to-image 生成的主流范式。
\textbf{Post-training with RL} 的目标是在预训练模型基础上，通过外部 reward signal（人类偏好、任务指标）进一步对齐生成结果，形式化为：

\[
\max_\theta \mathbb{E}_{x \sim \pi_\theta}[R(x)] - \beta \, \mathrm{KL}(\pi_\theta \| \pi_\mathrm{ref})
\]

最优解为 \(\pi^*(x) \propto \pi_\mathrm{ref}(x) \exp(R(x)/\beta)\)。

LLM 领域 PPO/GRPO 大获成功，但直接迁移到 diffusion model 存在根本障碍：\textbf{diffusion model 不是 likelihood-based 模型，无法直接计算 \(\log \pi_\theta(x_0)\)}，而 policy gradient 方法恰恰依赖这一项。现有工作（FlowGRPO 等）对此的处理方式不同，但缺乏系统性分析。

\hrule

\section{设计空间的三个维度}

本文将 RL fine-tuning 的设计空间解耦为三个正交维度，逐一控制变量分析：

\begin{center}
\begin{tabular}{ll}
\toprule
维度 & 候选方案 \\
\midrule
Policy-gradient objective & EPG / PEPG / PAR / GRPO \\
Likelihood estimator & Trajectory-based（FlowGRPO）/ ELBO-based \\
Sampling strategy & SDE / ODE \\
\bottomrule
\end{tabular}
\end{center}

\textbf{核心结论：Likelihood estimation 是主导因素，policy-gradient loss 的选择影响甚微。}

\hrule

\section{Likelihood Estimation：两种路线}

\subsection{Trajectory-based estimator（FlowGRPO）}

利用反向 SDE 的 Gaussian transition 逐步分解 likelihood：

\[
\log \pi_\theta(x_0) = \sum_{i=1}^N \log p_\theta(x_{t_{i-1}} | x_{t_i}), \quad
p_\theta(x_{t_{i-1}}|x_{t_i}) = \mathcal{N}\!\left(\mu_\theta,\; g_{t_i}^2(t_i - t_{i-1})I\right)
\]

\textbf{缺点：}

\begin{itemize}
  \item 必须使用 SDE sampler（ODE 下 \(p_\theta(x_{t_{i-1}}|x_{t_i})\) 退化为 delta，估计失效）
  \item 需要存储完整的采样轨迹（所有中间 \(x_t\)），内存开销与 NFE 数线性相关
  \item 训练时 sampler 与推理时不同，存在 train-inference mismatch
\end{itemize}

\subsection{ELBO-based estimator（本文）}

利用 flow matching 的训练目标直接近似 \(\log \pi_\theta(x_0)\)：

\[
\log \pi_\theta(x_0) \approx -\mathbb{E}_{t,\epsilon}\!\left[w(t)\|v_\theta(x_t, t) - v\|^2\right] + C_\mathrm{fw}(x_0)
\]

其中 \(v = \epsilon - x_0\)（flow model 的 conditional velocity），\(C_\mathrm{fw}\) 与 \(\theta\) 无关，不影响梯度。

\textbf{优点：}

\begin{itemize}
  \item 只需最终生成样本 \(x_0\)，\textbf{无需存储轨迹}，节省大量显存
  \item 与 sampler 完全解耦，可使用任意 black-box ODE solver
  \item 计算量远低于 trajectory 方案
\end{itemize}

\subsubsection{ELBO Weighting 的三种变体}

\begin{center}
\begin{tabular}{lll}
\toprule
名称 & \(w(t)\) & 备注 \\
\midrule
Path-KL & \(\frac{1-t}{t}\) & 来自 score-based 原始推导，实验中\textbf{训练不稳定} \\
Simple & \(1\) & 常数权重，简单有效 \\
Adaptive & \(\frac{t \cdot d \cdot \|v_\theta - v\|_2^2}{\mathrm{sg}(\|v_\theta - v\|_1)}\) & 数据依赖的自归一化，\textbf{默认推荐} \\
\bottomrule
\end{tabular}
\end{center}

\textbf{ELBO 的 Monte Carlo 估计方案：}

\begin{itemize}
  \item \emph{Single-timestep}：每次随机抽一个 \(t_i\)，估计一步，计算最廉价
  \item \emph{All-timestep}：对完整离散时间轴 \(\{1/N, \ldots, 1\}\) 求均值，方差更小
\end{itemize}

实验表明两者性能相当，\textbf{推荐 single-timestep} 以降低计算开销。

\hrule

\section{Policy-Gradient Objectives}

\subsection{EPG（Exact Policy Gradient）}

REINFORCE + group mean baseline，相比 GRPO 做了三处简化：

\[
\mathcal{L}_\mathrm{epg}(\theta) = \mathbb{E}_{x^i \sim \pi_{\theta_\mathrm{old}}}\!\left[\mathrm{sg}(\rho_\theta)(x^i) A^i_\mathrm{epg} \log \pi_\theta(x^i) - \beta\, \mathrm{kl}(x^i)\right]
\]

\begin{enumerate}
  \item \textbf{去掉 Clipping}：\(\epsilon\) 难以调优，在 diffusion 设置下更不稳定，实验证明去掉无损
  \item \textbf{去掉 Std Normalization}：连续 reward 分布下 std 差异小，normalization 引入 prompt-level bias
  \item \textbf{去掉 CFG}：CFG 使采样分布更尖锐/集中，与无 guidance 的训练分布不匹配，且双倍 NFE
\end{enumerate}

\subsection{PEPG（Proximal Exact Policy Gradient）}

在概率测度空间上做 proximal gradient descent，理论上等价于 trust-region 约束但无需 clipping heuristic：

\[
\mathcal{L}_\mathrm{pepg}(\theta) = \mathbb{E}\!\left[\left(A^i - \log \mathrm{sg}(\rho_\theta)(x^i)\right)\mathrm{sg}(\rho_\theta)(x^i)\log \pi_\theta(x^i) - \eta\, \mathrm{kl}(x^i)\right]
\]

迭代更新存在闭式解：

\[
\pi_{k+1}(x_0) \propto \left[\exp(R(x_0)/\beta)\,\pi_\mathrm{ref}(x_0)\right]^{\frac{\eta}{1+\eta}} \pi_k(x_0)^{\frac{1}{1+\eta}}
\]

当 \(k \to \infty\) 时收敛至目标 \(\pi^*\)。

\subsection{PAR（Proximal Advantage Regression）}

用 \(L_2\) 回归让 log-ratio 匹配 advantage，结构类似 GFlowNet trajectory balance：

\[
\mathcal{L}_\mathrm{par}(\theta) = \mathbb{E}\!\left[\frac{1}{2}\mathrm{sg}(\rho_\theta)(x^i)\left\|A^i - \log \rho_\theta(x^i)\right\|^2 - \eta\, \mathrm{kl}(x^i)\right]
\]

\textbf{Theorem 3.1}：EPG、PEPG、PAR 三者共享相同的最优解 \(\pi^*(x) \propto \pi_\mathrm{ref}(x)\exp(R(x)/\beta)\)。

\hrule

\section{实验结果}

\subsection{主要效率对比（GenEval 0.95）}

\begin{center}
\begin{tabular}{ll}
\toprule
方法 & GPU Hours (8×H100) \\
\midrule
FlowGRPO (w/ CFG) & 422.9 \\
DiffusionNFT & 340.8 \\
Ours (ELBO, SDE) & 185.4 \\
\textbf{Ours (ELBO, ODE)} & \textbf{90.4} \\
\bottomrule
\end{tabular}
\end{center}

\begin{itemize}
  \item ELBO vs Trajectory：\textbf{4.68× 加速}（prompt efficiency）
  \item ODE vs SDE（ELBO 下）：性能持平，但 ODE 仅需 10 steps（vs SDE 40 steps），\textbf{额外 \textasciitilde{}2× 加速}
\end{itemize}

\subsection{多 Benchmark 对比（Tab. 2）}

\begin{center}
\begin{tabular}{llllll}
\toprule
方法 & GenEval & OCR & PickScore & HPSv2.1 & Aesthetic \\
\midrule
FlowGRPO & 0.95 & 0.92 & 22.51 & 0.274 & 5.32 \\
DiffusionNFT & 0.95 & 0.93 & 22.88 & 0.289 & 5.25 \\
\textbf{Ours} & \textbf{0.96} & \textbf{0.94} & \textbf{22.93} & \textbf{0.289} & \textbf{5.33} \\
\bottomrule
\end{tabular}
\end{center}

\subsection{Policy-gradient loss 的影响（Tab. 1，ELBO+ODE 下）}

\begin{center}
\begin{tabular}{lll}
\toprule
Loss & GenEval & PickScore \\
\midrule
EPG & 0.96 & 22.77 \\
PEPG & 0.96 & 22.85 \\
PAR & 0.96 & 22.97 \\
GRPO & 0.94 & 22.45 \\
\bottomrule
\end{tabular}
\end{center}

差异极小，验证了 \textbf{loss 选择不是主导因素}。

\hrule

\section{与相关工作的关系}

\begin{verbatim}
FlowGRPO   →  GRPO loss  + Trajectory estimator + SDE  （重、慢）
AWM        →  GRPO loss  + ELBO (simple, 1-sample) + ODE  
DiffusionNFT → NFT-style + ELBO (adaptive) + ODE（多阶段训练，SOTA）
本文        →  极简 loss (PEPG/PAR/EPG) + ELBO (adaptive) + ODE（单阶段、更快）
\end{verbatim}

AWM 本质上已经用了 ELBO+ODE，但沿用了 GRPO 的 loss 工程；本文进一步剥离了 loss 设计的影响，证明极简 EPG 即可达到同等效果。

\hrule

\section{实践建议（来自论文）}

\begin{itemize}
  \item \textbf{默认配置}：PEPG loss + Adaptive ELBO + ODE sampler (10 steps) + 无 CFG
  \item ELBO weighting 避免使用 Path-KL（训练不稳定）
  \item Single-timestep Monte Carlo 估计足够，无需全轨迹
  \item Clipping、std normalization、CFG 均可去掉
  \item KL divergence 估计采用 simple weighting (\(w(t)=1\))，实现简单且有效
\end{itemize}
